\begin{abstract}
We study semi-supervised learning on a clinical-trial corpus in which only 157 of
1,171 training trials carry stopping-reason labels. \textbf{Part~1} predicts five
binary stopping reasons (Safety, Funding, Efficacy, Recruitment, Other) from 150-d
reason embeddings. Our final model couples an unsupervised diagonal-covariance
Gaussian mixture (30 components, fit on all 1,561 trials without labels) to a
class-balanced multinomial logistic regression blended with a cosine $k$NN, and
uses the mixture posterior as a prior inside two rounds of confidence-thresholded
self-training. It reaches sample-F1 0.893 and macro-F1 0.714 on the validation set
(0.914 / 0.814 under repeated cross-validation) and 0.894 on the Kaggle test set.
An ablation shows that eligibility criteria, disease text, Morgan fingerprints and
start year all fail to improve on the embeddings alone. \textbf{Part~2} predicts
per-phase trial success. A stacked ensemble of eight field-specific base models,
combined by a phase-aware logistic meta-model and trained with exponential recency
weights, cuts mean validation log loss from 0.591 to 0.556 and raises mean PR-AUC
from 0.796 to 0.814 relative to a criteria-only per-phase baseline. Error analysis
keyed on the Part~1 reasons shows that failures concentrate in Phase~III and in
trials with no recorded stopping reason, and that the improved model reduces those
error rates most.
\end{abstract}
